% This document was generated by an LLM.

\documentclass{essay}

\title{The Infinite Nested Radical as a Limit}
\author{}
\date{}

\begin{document}

\maketitle

\section{Introduction}

Consider the infinite nested radical
\[
\sqrt{2+\sqrt{2+\sqrt{2+\cdots}}}.
\]
A familiar informal argument assigns a value \(L\) to the entire expression and observes that, after removing the outermost square root, the same infinite expression appears again. One therefore writes
\[
L=\sqrt{2+L},
\]
and hence
\[
L^2=L+2.
\]
Solving the quadratic equation gives
\[
L\in\{-1,2\}.
\]
Since a square root is nonnegative, one concludes that \(L=2\).

The numerical conclusion is correct, but the argument is logically incomplete. It presupposes that the infinite radical has a finite value before the existence of such a value has been established. In rigorous analysis, an infinite expression involving an ellipsis is not automatically a number. It must first be defined by means of a limiting process, and the relevant limit must then be shown to exist.

The purpose of this note is to formulate the infinite radical precisely and prove that its value is \(2\). The proof illustrates several basic ideas of real analysis: sequences, induction, monotonicity, boundedness, completeness of the real numbers, continuity, and fixed points.

\section{Defining the Infinite Radical}

For each positive integer \(n\), define a finite nested radical by the recurrence
\[
a_1=\sqrt{2},
\qquad
a_{n+1}=\sqrt{2+a_n}.
\]
Thus
\[
a_1=\sqrt{2},
\]
\[
a_2=\sqrt{2+\sqrt{2}},
\]
\[
a_3=\sqrt{2+\sqrt{2+\sqrt{2}}},
\]
and so forth.

The infinite nested radical is then defined, if the limit exists, by
\[
\sqrt{2+\sqrt{2+\sqrt{2+\cdots}}}
:=
\lim_{n\to\infty} a_n.
\]

Accordingly, there are two distinct questions:

\begin{enumerate}
    \item Does the sequence \((a_n)\) converge?
    \item If it converges, what is its limit?
\end{enumerate}

The second question is relatively easy. The first is the essential point that the informal argument omits.

\section{Convergence of the Sequence}

We shall prove that \((a_n)\) is increasing and bounded above. The monotone convergence theorem for sequences will then imply that \((a_n)\) converges.

\begin{theorem}[Monotone Convergence Theorem]
Every monotone increasing sequence of real numbers that is bounded above converges in \(\mathbb{R}\).
\end{theorem}

This theorem is ultimately a consequence of the completeness of the real numbers: every nonempty subset of \(\mathbb{R}\) that is bounded above has a least upper bound.

We begin with boundedness.

\begin{proposition}
For every \(n\in\mathbb{N}\),
\[
0<a_n<2.
\]
\end{proposition}

\begin{proof}
We argue by induction.

For \(n=1\),
\[
a_1=\sqrt{2},
\]
so
\[
0<a_1<2.
\]

Now suppose that
\[
0<a_n<2
\]
for some \(n\geq 1\). Then
\[
2<2+a_n<4.
\]
Since the square-root function is strictly increasing on \([0,\infty)\),
\[
\sqrt{2}<\sqrt{2+a_n}<2.
\]
Therefore
\[
0<a_{n+1}<2.
\]

By induction, the inequality holds for all \(n\in\mathbb{N}\).
\end{proof}

Thus the sequence is bounded above by \(2\). We next establish monotonicity.

\begin{proposition}
The sequence \((a_n)\) is strictly increasing:
\[
a_{n+1}>a_n
\qquad
\text{for every } n\in\mathbb{N}.
\]
\end{proposition}

\begin{proof}
Fix \(n\in\mathbb{N}\). From the previous proposition,
\[
0<a_n<2.
\]
Because both \(a_n\) and \(a_{n+1}\) are positive, the inequality
\[
a_{n+1}>a_n
\]
is equivalent to
\[
a_{n+1}^2>a_n^2.
\]
Using the recurrence relation
\[
a_{n+1}^2=2+a_n,
\]
we therefore need to prove
\[
2+a_n>a_n^2.
\]
Equivalently,
\[
a_n^2-a_n-2<0.
\]
Factoring gives
\[
(a_n-2)(a_n+1)<0.
\]
Since
\[
a_n-2<0
\qquad\text{and}\qquad
a_n+1>0,
\]
their product is negative. Hence
\[
a_{n+1}>a_n.
\]
\end{proof}

Combining the two propositions, we obtain
\[
0<a_1<a_2<a_3<\cdots<2.
\]
Therefore \((a_n)\) is increasing and bounded above.

\begin{theorem}
The sequence \((a_n)\) converges.
\end{theorem}

\begin{proof}
The sequence is monotone increasing and bounded above. Hence, by the monotone convergence theorem, there exists some \(L\in\mathbb{R}\) such that
\[
\lim_{n\to\infty} a_n=L.
\]
\end{proof}

At this stage the existence of a finite limit has been established. Only now is it legitimate to determine the limit by using the self-similar structure of the recurrence.

\section{Identification of the Limit}

The recurrence relation is
\[
a_{n+1}=\sqrt{2+a_n}.
\]
Since
\[
a_n\to L,
\]
the shifted sequence \((a_{n+1})\) has the same limit:
\[
a_{n+1}\to L.
\]

The function
\[
f(x)=\sqrt{2+x}
\]
is continuous on \([-2,\infty)\). Therefore continuity permits the interchange of \(f\) with the limit:
\[
\lim_{n\to\infty} f(a_n)
=
f\left(\lim_{n\to\infty} a_n\right).
\]
Consequently,
\[
L
=
\lim_{n\to\infty} a_{n+1}
=
\lim_{n\to\infty}\sqrt{2+a_n}
=
\sqrt{2+L}.
\]

Squaring both sides yields
\[
L^2=L+2,
\]
so
\[
L^2-L-2=0.
\]
Thus
\[
(L-2)(L+1)=0,
\]
and hence
\[
L=2
\qquad\text{or}\qquad
L=-1.
\]

However, every \(a_n\) is positive. Since a limit of a convergent sequence of nonnegative real numbers is nonnegative,
\[
L\geq 0.
\]
Therefore \(L=-1\) is impossible, and we conclude that
\[
\boxed{L=2}.
\]

Hence
\[
\boxed{
\sqrt{2+\sqrt{2+\sqrt{2+\cdots}}}=2
}.
\]

\section{Why the Informal Argument Is Incomplete}

The equation
\[
L=\sqrt{2+L}
\]
does not by itself prove that the nested radical converges.

Indeed, suppose in general that a sequence is generated by iteration
\[
x_{n+1}=f(x_n).
\]
If \(x_n\to L\) and \(f\) is continuous at \(L\), then necessarily
\[
L=f(L).
\]
Thus every limit of the iteration must be a fixed point of \(f\).

The converse, however, is false. The existence of a fixed point does not imply that a particular iterative sequence converges to it, or even that the sequence converges at all.

For example, consider
\[
x_{n+1}=1-x_n.
\]
The function
\[
f(x)=1-x
\]
has the unique fixed point
\[
L=\frac12.
\]
Nevertheless, if \(x_1=0\), then
\[
0,1,0,1,0,1,\ldots
\]
is obtained, and the sequence does not converge.

Therefore solving the fixed-point equation
\[
L=f(L)
\]
answers only the conditional question:

\begin{quote}
If the iteration converges, what values are possible for its limit?
\end{quote}

It does not answer the prior question of whether convergence occurs.

For the nested radical, the rigorous proof must therefore proceed in the order
\[
\text{definition}
\longrightarrow
\text{convergence}
\longrightarrow
\text{fixed-point equation}
\longrightarrow
\text{identification of the limit}.
\]

\section{A Fixed-Point Interpretation}

The recurrence
\[
a_{n+1}=\sqrt{2+a_n}
\]
may be viewed as an iteration of the function
\[
f(x)=\sqrt{2+x}.
\]
The number \(2\) is a fixed point because
\[
f(2)=\sqrt{4}=2.
\]

The derivative of \(f\) is
\[
f'(x)=\frac{1}{2\sqrt{2+x}}.
\]
At the fixed point,
\[
f'(2)=\frac14.
\]
Since
\[
|f'(2)|<1,
\]
the fixed point is locally attracting. Informally, points sufficiently close to \(2\) are moved closer to \(2\) under repeated application of \(f\).

This can also be made quantitative. Let
\[
e_n=2-a_n.
\]
Since \(a_n<2\), we have \(e_n>0\). Then
\begin{align*}
e_{n+1}
&=
2-\sqrt{2+a_n}\\
&=
\frac{4-(2+a_n)}
     {2+\sqrt{2+a_n}}\\
&=
\frac{2-a_n}
     {2+a_{n+1}}\\
&=
\frac{e_n}{2+a_{n+1}}.
\end{align*}
Because \(a_{n+1}>0\),
\[
2+a_{n+1}>2,
\]
and therefore
\[
0<e_{n+1}<\frac12 e_n.
\]
Thus the error decreases at least geometrically.

In fact, since \(a_{n+1}\to 2\),
\[
\frac{e_{n+1}}{e_n}
=
\frac{1}{2+a_{n+1}}
\longrightarrow
\frac14.
\]
Hence asymptotically
\[
e_{n+1}\sim \frac14 e_n.
\]
This agrees with the derivative calculation
\[
f'(2)=\frac14.
\]

Thus the fixed-point viewpoint not only explains the value of the limit but also describes the rate at which the finite radicals approach it.

\section{Completeness and the Role of Real Analysis}

The most substantial theorem used in the convergence proof is the monotone convergence theorem. It is worth emphasizing that this result depends on a specifically real-analytic property of \(\mathbb{R}\): completeness.

Let
\[
S=\{a_n:n\in\mathbb{N}\}.
\]
We have shown that \(S\) is nonempty and bounded above. Completeness guarantees the existence of
\[
L=\sup S.
\]
Since the sequence is increasing, its terms approach this least upper bound.

Indeed, let \(\varepsilon>0\). Because \(L-\varepsilon\) is not an upper bound for \(S\), there exists \(N\in\mathbb{N}\) such that
\[
a_N>L-\varepsilon.
\]
Since the sequence is increasing,
\[
a_n\geq a_N>L-\varepsilon
\qquad
\text{for all }n\geq N.
\]
On the other hand, \(L\) is an upper bound, so
\[
a_n\leq L.
\]
Therefore, for all \(n\geq N\),
\[
L-\varepsilon<a_n\leq L,
\]
which implies
\[
|a_n-L|<\varepsilon.
\]
Hence
\[
a_n\to L.
\]

This is an explicit \(\varepsilon\)-proof of the monotone convergence theorem in this particular case.

The argument shows where the existence of the limit really comes from. It does not come from algebraic self-similarity. It comes from the order structure and completeness of the real numbers.

\section{Conclusion}

The infinite nested radical
\[
\sqrt{2+\sqrt{2+\sqrt{2+\cdots}}}
\]
is rigorously understood as the limit of the recursively defined sequence
\[
a_1=\sqrt2,
\qquad
a_{n+1}=\sqrt{2+a_n}.
\]

One first proves that
\[
0<a_n<2
\]
for every \(n\), and that
\[
a_{n+1}>a_n.
\]
Hence the sequence is increasing and bounded above. By completeness of the real numbers, it converges.

Only after convergence has been established may one let
\[
L=\lim_{n\to\infty}a_n
\]
and pass to the limit in the recurrence. Continuity of the square-root function gives
\[
L=\sqrt{2+L},
\]
from which
\[
L=2.
\]

The principal methodological lesson is broader than this particular radical. Whenever an ``infinite expression'' is written using an ellipsis, one should ask what limiting process the notation represents and whether that limit exists. Algebraic manipulation of a hypothetical limit may identify possible values, but it cannot substitute for a proof of convergence.

\end{document}
